%% ch12_cybersecurity.tex
%% Chapter 12: Cybersecurity for Adaptive Protection Systems
%% ============================================================
\chapter{Cybersecurity for Adaptive Protection Systems}
\label{ch:cybersecurity}

\section{Introduction}
\label{sec:ch12_intro}

The shift from fixed relay settings to real-time adaptive protection
introduces a new attack surface: the communication and computation
infrastructure that carries the data on which adaptation depends.
An adversary who can inject false current measurements, spoof GOOSE
messages, or manipulate the PINN classifier input has the ability to
cause the SAMBP system to either \emph{refuse to trip} (security failure)
or \emph{trip incorrectly} (reliability failure)---outcomes that are
equivalent in consequence to the physical protection failures described
in Chapter~\ref{ch:ibr_char}.

This chapter addresses three cybersecurity dimensions of the SAMBP framework:

\begin{enumerate}
\item \textbf{GOOSE message integrity (Section~\ref{sec:ch12_goose}):}
      HMAC-SHA256 authentication countermeasure for IEC~61850 GOOSE
      spoofing and replay attacks. This contribution (Contribution C17)
      was developed and validated in Chapter~\ref{ch:centralised}; this
      chapter provides the formal security analysis and threat model.

\item \textbf{Adversarial PINN robustness (Section~\ref{sec:ch12_adversarial}):}
      Analysis and mitigation of CT-injection attacks targeting the PINN
      fault classifier (Contribution C25 [F]).

\item \textbf{Coordinated attack scenarios (Section~\ref{sec:ch12_coordinated}):}
      Formal threat model for coordinated cyberattacks targeting the
      closed-loop adaptation protocol (CUSUM$\to$settings$\to$GOOSE).
\end{enumerate}

\section{IEC 61850 GOOSE Security Architecture}
\label{sec:ch12_goose}

\subsection{GOOSE Protocol and Attack Surface}

IEC~61850 GOOSE (Generic Object Oriented Substation Event) messages are
multicast UDP frames transmitted on the substation LAN at rates up to
1~ms per message. They carry protection trip signals, interlocking status,
and breaker position indications. GOOSE has no built-in authentication:
any host on the substation VLAN can inject a GOOSE message with any
dataset content.

The primary attacks are~\citep{IEC62351}:
\begin{description}
\item[Spoofing:] Injecting a fake GOOSE trip signal to cause an
    unintended circuit breaker operation.
\item[Replay:] Replaying a captured genuine trip message at a later
    time (when the system is healthy) to cause a spurious outage.
\item[Suppression:] Blocking genuine GOOSE messages so that a real
    fault is not cleared.
\end{description}

\subsection{HMAC-SHA256 Countermeasure (Contribution C17)}

The SAMBP GOOSE security layer (implemented in Chapter~\ref{ch:centralised})
appends a 32-byte HMAC-SHA256 tag to each GOOSE application layer dataset:
\begin{equation}
  \mathrm{tag}_t = \mathrm{HMAC}_k\bigl(\mathrm{DataSet}_t \,\|\, t_\mathrm{UTC} \,\|\, \mathrm{SqNum}\bigr)
  \label{eq:ch12_hmac}
\end{equation}
where $k$ is a 256-bit shared key established at commissioning,
$t_\mathrm{UTC}$ is the microsecond-precision IEEE~1588 PTP timestamp,
and $\mathrm{SqNum}$ is the GOOSE sequence number. The receiver verifies
the tag before acting on any GOOSE dataset.

This construction prevents:
\begin{itemize}
\item Spoofing: An adversary without $k$ cannot forge a valid tag.
\item Replay: The timestamp and sequence number uniqueness condition
      rejects any message with $|t_\mathrm{UTC} - t_\mathrm{now}| > \delta_\mathrm{replay} = 5$~ms.
\end{itemize}

\subsection{Security and Performance Results}

From the WAPC validation of Chapter~\ref{ch:centralised}:

\begin{table}[htbp]
\caption{HMAC-SHA256 GOOSE security results}
\label{tab:ch12_goose_results}
\begin{tabularx}{\textwidth}{@{}lXXl@{}}
\toprule
Metric & Before HMAC & After HMAC & Status \\
\midrule
GOOSE spoofing success rate & 12.3\% & $< 0.001\%$ & [P] \\
GOOSE replay success rate & 8.7\% & $0.0\%$ & [P] \\
GOOSE latency overhead & 0.00~ms & 0.72~ms & [P] \\
P99 trip time (with HMAC) & & $< 10$~ms & [P] \\
Risk reduction (risk matrix) & -- & $-54.9\%$ & [P] \\
\bottomrule
\end{tabularx}
\end{table}

The 0.72~ms latency overhead is within the IEC~61850 real-time requirement
($< 4$~ms for class P2 messages). The 54.9\% risk reduction places the
system in the ``acceptable'' zone of the IEC~62351 risk matrix.

\subsection{Remaining Vulnerabilities}

HMAC addresses spoofing and replay but does not mitigate:
\begin{enumerate}
\item \textbf{Flooding (DoS):} An adversary sending high-rate GOOSE
      traffic can exhaust the substation LAN bandwidth. Mitigation:
      switch-level rate limiting and VLAN segmentation (planned in G5,
      Section~\ref{sec:ch13_future}).
\item \textbf{Key compromise:} If the shared key $k$ is extracted from
      a compromised IED, all HMAC protection is lost. Mitigation: automatic
      key rotation per IEC~62351-8 (planned future work).
\end{enumerate}

\section{Adversarial PINN Robustness (Contribution C25)}
\label{sec:ch12_adversarial}

\subsection{Threat Model: CT-Injection Attack}

A CT-injection (current transformer injection) attack inserts a small
crafted signal $\bm{\delta}$ into the CT secondary current:
\begin{equation}
  \tilde{\mathbf{x}} = \mathbf{x} + \bm{\delta}, \quad
  \|\bm{\delta}\|_\infty \leq \varepsilon
  \label{eq:ch12_ct_attack}
\end{equation}
where $\mathbf{x}$ is the true current measurement and $\varepsilon$
bounds the perturbation magnitude. The attacker's goal is to cause
the PINN classifier to misclassify a genuine fault as a normal operating
condition (causing it not to trip) or to misclassify normal operation
as a fault (causing a false trip).

This is the power system equivalent of the adversarial example
attack~\citep{Goodfellow2015} in computer vision. The key question is:
does the physics constraint in the PINN loss function provide any inherent
robustness against such attacks?

\subsection{Theoretical Analysis: Physics Constraints as Robustness Mechanism}

\textbf{Hypothesis:} The physics-constrained PINN is more robust to
adversarial perturbations than an unconstrained neural network, because
any perturbation $\bm{\delta}$ that causes misclassification must also
violate the embedded KVL/KCL constraints (which the loss function penalises
strongly).

\textbf{Formal statement:} Let $f: \mathbb{R}^d \to [0,1]^c$ be the
PINN classifier, where $c$ is the number of fault classes. The
\emph{adversarial margin} is:
\begin{equation}
  m_\mathrm{adv}(\mathbf{x}) = \min_{j \neq y^\ast}
    \left[\min_{\|\bm{\delta}\|_\infty \leq \varepsilon}
    f_{y^\ast}(\mathbf{x} + \bm{\delta}) - f_j(\mathbf{x} + \bm{\delta})\right]
  \label{eq:ch12_adv_margin}
\end{equation}
where $y^\ast$ is the true class. A positive adversarial margin means the
classifier is robust to $\varepsilon$-perturbations.

\textbf{Claim:} For the PINN with physics penalty weight
$\lambda_\mathrm{phys} \gg \lambda_\mathrm{CE}$ (cross-entropy weight),
the adversarial margin is bounded below by a function of
$\lambda_\mathrm{phys}$ and the Lipschitz constant of the physics residual:
\begin{equation}
  m_\mathrm{adv}(\mathbf{x}) \geq \frac{\lambda_\mathrm{phys}}{L_\mathrm{phys}}\,
    g(\mathbf{x}, \varepsilon)
  \label{eq:ch12_margin_bound}
\end{equation}
where $g(\mathbf{x}, \varepsilon)$ is the physics constraint violation
of a perturbed input (to be derived).

\subsection{Planned Validation}
\label{sec:ch12_adversarial_validation}

\begin{tcolorbox}[colback=orange!5,colframe=orange!60!black,
  fonttitle=\bfseries\small,
  title={C25 --- Adversarial PINN Robustness (Pending)},
  breakable]
\small
\textbf{Study plan:}
\begin{enumerate}
\item Generate adversarial examples using:
      \begin{itemize}
      \item FGSM (Fast Gradient Sign Method): $\bm{\delta} = \varepsilon\,\mathrm{sgn}(\nabla_\mathbf{x} L)$
      \item PGD (Projected Gradient Descent) with 20 steps
      \item CT-injection physical model: restrict $\bm{\delta}$ to
            secondary current channels only
      \end{itemize}
\item Evaluate adversarial accuracy of:
      (a) unconstrained NN baseline, and
      (b) PINN with $\lambda_\mathrm{phys} \in [1, 100]$
      under $\varepsilon \in [0.01, 0.20]$~pu.
\item Adversarial retraining: augment training data with FGSM examples
      at $\varepsilon = 0.05$~pu, re-evaluate clean and adversarial accuracy.
\end{enumerate}

\textbf{Expected results:}
\begin{itemize}
\item PINN with $\lambda_\mathrm{phys} = 10$ is expected to retain $> 80\%$
      accuracy under $\varepsilon = 0.05$~pu perturbations, vs.\ $< 50\%$
      for the unconstrained baseline (hypothesis based on physics-constraint
      robustness theory).
\item Adversarial retraining expected to further improve adversarial
      accuracy to $> 90\%$ with $< 2$~pp degradation in clean accuracy.
\end{itemize}

\textbf{Status:} Threat model and study plan complete.
FGSM/PGD implementation pending.
\end{tcolorbox}

\section{Coordinated Attack Threat Model}
\label{sec:ch12_coordinated}

\subsection{Multi-Vector Attack on the Closed-Loop Protocol}

The most sophisticated attack targets the closed-loop CUSUM$\to$settings$\to$GOOSE
protocol of Chapter~\ref{ch:centralised} by simultaneously:
\begin{enumerate}
\item Injecting a false $\kibr$ estimate into the EKF via CT manipulation.
\item Triggering a false CUSUM alarm, causing the WAPC to initiate
      settings re-optimisation.
\item Spoofing the GOOSE confirmation message to cause the relay to
      accept the new (malicious) settings without physical verification.
\end{enumerate}

This is a three-layer attack that bypasses each individual countermeasure
(HMAC on GOOSE, confidence gate $\gamma$ on EKF, CUSUM threshold) by
exploiting the interaction between layers.

\subsection{Defence-in-Depth Framework}

The SAMBP multi-layer defence against coordinated attacks:
\begin{enumerate}
\item \textbf{HMAC on GOOSE} (Layer 1): Prevents step 3 above without key compromise.
\item \textbf{Confidence gate $\gamma \geq 0.70$} (Layer 2): A CT injection
      that is large enough to shift $\kibr$ by a meaningful amount will
      also degrade $\gamma_J$ (Jacobian condition) and $\gamma_P$
      (physics plausibility), triggering the adaptation block.
\item \textbf{Bounded update law} (Layer 3): Even if a false $\kibr$
      estimate passes the confidence gate, the bounded update law
      limits the maximum settings change to 10\% per cycle, requiring
      an attacker to sustain the injection for $> 10$ cycles to cause
      significant harm.
\item \textbf{CUSUM rate-of-change alarm} (Layer 4): A CUSUM detector on
      the settings trajectory (not just the measurement trajectory)
      will trigger an anomaly alert if settings change faster than
      physically plausible for normal fleet dynamics.
\end{enumerate}

\subsection{Quantitative Defence Analysis}

\begin{tcolorbox}[colback=orange!5,colframe=orange!60!black,
  fonttitle=\bfseries\small,
  title={Coordinated Attack Quantitative Analysis (Pending)},
  breakable]
\small
\textbf{Study plan:} Monte Carlo simulation of 1000 coordinated attack
scenarios across:
\begin{itemize}
\item CT injection magnitudes $\varepsilon \in [0.01, 0.30]$~pu
\item Attack durations $T_\mathrm{attack} \in [1, 60]$~s
\item Network topologies: radial, meshed, microgrid
\end{itemize}

\textbf{Metrics:}
\begin{itemize}
\item False trip probability under attack
\item False restrain probability under attack
\item Mean time before the attack is detected by Layer~4 CUSUM
\end{itemize}

\textbf{Status:} Framework designed. Monte Carlo engine pending implementation.
\end{tcolorbox}

\section{Standards Compliance Summary}
\label{sec:ch12_standards}

\begin{table}[htbp]
\caption{Cybersecurity standards compliance status}
\label{tab:ch12_standards}
\begin{tabularx}{\textwidth}{@{}lp{5cm}Xl@{}}
\toprule
Standard & Requirement & SAMBP implementation & Status \\
\midrule
IEC~62351-6 & GOOSE authentication & HMAC-SHA256 per Eq.~\eqref{eq:ch12_hmac} & [P] \\
IEC~62351-8 & Key management & Manual commissioning (auto-rotation planned) & Partial \\
NERC CIP-005 & Network access control & VLAN segmentation (designed) & [F] \\
NERC CIP-007 & Security event monitoring & CUSUM anomaly alarm (Layer~4) & Partial \\
IEEE~1686 & IED cyber security & SIL proxy validation completed & [P] \\
\bottomrule
\end{tabularx}
\end{table}

\section{Summary}
\label{sec:ch12_summary}

This chapter has established the cybersecurity architecture for the SAMBP
adaptive protection system. The key results are:

\begin{description}
\item[C17 (validated):] HMAC-SHA256 reduces GOOSE attack risk by 54.9\%
      at 0.72~ms latency overhead, meeting IEC~61850 Class~P2 requirements.
\item[C25 (framework complete):] The theoretical basis for adversarial
      PINN robustness is established. The physics-constraint mechanism
      is hypothesised to provide $> 30\%$ adversarial accuracy improvement
      over unconstrained baselines. Experimental validation is pending.
\end{description}

The four-layer defence-in-depth framework (HMAC $\to$ confidence gate
$\to$ bounded update $\to$ CUSUM trajectory monitoring) provides a
comprehensive response to coordinated multi-vector attacks without
requiring any single layer to be unbreakable. The weakest link is the
absence of automatic key rotation (IEC~62351-8), which is identified as
a priority future implementation item.
